\documentclass[12pt,a4paper]{article}

%% ── Packages ──────────────────────────────────────────────────────────────
\usepackage[utf8]{inputenc}
\usepackage[T1]{fontenc}
\usepackage{lmodern}
\usepackage[margin=2.5cm]{geometry}
\usepackage{amsmath,amssymb}
\usepackage{graphicx}
\usepackage{booktabs}
\usepackage{array}
\usepackage{multirow}
\usepackage{hyperref}
\usepackage{xcolor}
\usepackage{listings}
\usepackage{caption}
\usepackage{subcaption}
\usepackage{setspace}
\usepackage{titlesec}
\usepackage[numbers,sort&compress]{natbib}
\usepackage{url}

\hypersetup{
  colorlinks=true,
  linkcolor=blue!70!black,
  citecolor=blue!70!black,
  urlcolor=blue!70!black
}

\lstset{
  basicstyle=\ttfamily\footnotesize,
  breaklines=true,
  frame=single,
  backgroundcolor=\color{gray!8},
  keywordstyle=\color{blue!80!black}\bfseries,
  commentstyle=\color{green!50!black},
  stringstyle=\color{red!70!black},
  numbers=left,
  numberstyle=\tiny\color{gray},
  xleftmargin=1.5em
}

\onehalfspacing
\setlength{\parskip}{6pt}
\setlength{\parindent}{0pt}

%% ── Title ─────────────────────────────────────────────────────────────────
\title{%
  \vspace{-1cm}
  {\Large \textbf{Explainable Lightweight Intrusion Detection System}} \\[6pt]
  {\large for IoT Networks Using Ensemble Learning}
}

\author{
  Abdelrahman Fakhry \\
  \textit{Egypt-Japan University of Science and Technology (E-JUST)} \\
  \texttt{abdelrahman.fakhry@ejust.edu.eg}
}

\date{\today}

\begin{document}

\maketitle
\thispagestyle{empty}

%% ── Abstract ──────────────────────────────────────────────────────────────
\begin{abstract}
The proliferation of Internet of Things (IoT) devices has dramatically expanded the
network attack surface, making effective intrusion detection a critical security
requirement. Traditional signature-based Intrusion Detection Systems (IDS) fail to
identify zero-day and novel attack patterns, while purely black-box machine learning
approaches lack the transparency required for operational trust and forensic analysis.
This paper presents \textbf{XAI-IDS}, an Explainable Lightweight Intrusion Detection
System designed for IoT network environments. The system combines a Random Forest
ensemble classifier trained on the NF-UNSW-NB15-v3 NetFlow dataset with
SHapley Additive exPlanations (SHAP) to provide both global feature-level insights
and local per-prediction reasoning in human-readable natural language.
Evaluated on 2,365,424 network flows spanning nine attack categories, the binary
classifier achieves \textbf{99.99\% accuracy}, \textbf{99.94\% F1-score},
\textbf{0.02\% false negative rate}, and an AUC-ROC of 1.00.
An edge-optimised variant (25 trees, depth-8) retains 99.91\% F1 while achieving
\textbf{11,055 flows per second} throughput at \textbf{0.25~MB model size},
making it suitable for deployment on resource-constrained IoT gateways.
The system is delivered as a real-time Streamlit dashboard and a RESTful FastAPI
inference endpoint, enabling operational deployment without specialist ML expertise.
\end{abstract}

\textbf{Keywords:} Intrusion Detection System, IoT Security, Explainable AI, Random
Forest, SHAP, Network Anomaly Detection, Edge Computing.

\newpage
\tableofcontents
\newpage

%% ── 1. Introduction ───────────────────────────────────────────────────────
\section{Introduction}
\label{sec:intro}

The Internet of Things has transformed modern infrastructure, connecting billions of
devices ranging from industrial sensors and medical equipment to smart home appliances.
By 2025, the number of active IoT endpoints is projected to exceed 75 billion
\cite{ferrag2020deep}. This rapid expansion introduces significant security
vulnerabilities: IoT devices often operate with limited computational resources,
infrequent firmware updates, and weak authentication mechanisms, making them
attractive targets for adversaries.

Intrusion Detection Systems are a foundational component of network security, tasked
with identifying malicious traffic and alerting operators. Conventional IDS solutions
fall into two broad categories: \emph{signature-based systems}, which match known
attack patterns from a curated database, and \emph{anomaly-based systems}, which
learn a baseline of normal behaviour and flag deviations
\cite{garcia2009anomaly,liao2013iot}. While signature-based approaches achieve high
precision for known threats, they are inherently incapable of detecting zero-day
exploits. Anomaly-based systems powered by machine learning partially address this
limitation but introduce a new challenge: opacity.

A critical shortcoming of modern machine learning models deployed as IDS is the
lack of \emph{explainability}. When a deep neural network or a large ensemble flags
a connection as malicious, security analysts receive a binary verdict but no
justification. This opacity undermines operational trust, hampers forensic
investigation, impedes regulatory compliance (e.g., GDPR Article 22), and prevents
analysts from identifying whether a high-confidence decision is correct or an artefact
of overfitting \cite{doshi2017towards,gunning2019xai}.

The field of Explainable Artificial Intelligence (XAI) addresses this gap by
providing techniques that reveal the internal reasoning of otherwise black-box
models \cite{molnar2022}. Among these, SHAP (SHapley Additive exPlanations)
\cite{lundberg2017shap} offers a theoretically grounded, game-theoretic framework
for attributing each model prediction to its contributing features. SHAP has emerged
as the de facto standard for post-hoc model explanation due to its consistency,
local accuracy, and compatibility with tree-based ensemble methods.

Simultaneously, IoT deployments impose strict computational constraints. An IDS
that requires a GPU server or multi-gigabyte model weights is impractical for edge
gateways such as Raspberry Pi class devices operating at the network periphery.
There is therefore a pressing need for IDS solutions that are simultaneously
\emph{accurate}, \emph{explainable}, and \emph{lightweight}.

This paper makes the following contributions:

\begin{enumerate}
  \item \textbf{A complete XAI-IDS pipeline}: end-to-end preprocessing, Random
        Forest binary and multiclass classification, and SHAP-based explainability
        on the NF-UNSW-NB15-v3 dataset.
  \item \textbf{Natural language explanation generation}: automatic conversion of
        SHAP attribution values into human-readable reasoning statements, enabling
        non-specialist operators to understand alert causes.
  \item \textbf{Edge-optimised model}: a 0.25~MB, 25-tree variant achieving
        11,055 flows/sec at 0.09~ms per-sample latency, validated against
        Raspberry Pi 4 class deployment targets.
  \item \textbf{Operational dashboard and API}: a Streamlit real-time monitoring
        interface and a FastAPI REST endpoint for integration into existing
        security operations workflows.
  \item \textbf{Comprehensive evaluation}: binary and multiclass metrics including
        false negative rate analysis, cross-attack-class performance, and
        edge vs. full-model benchmarking.
\end{enumerate}

The remainder of this paper is structured as follows.
Section~\ref{sec:related} surveys related work.
Section~\ref{sec:dataset} describes the dataset and preprocessing.
Section~\ref{sec:methodology} presents the system architecture and methodology.
Section~\ref{sec:results} reports experimental results.
Section~\ref{sec:discussion} discusses findings and limitations.
Section~\ref{sec:conclusion} concludes and outlines future directions.

%% ── 2. Related Work ───────────────────────────────────────────────────────
\section{Related Work}
\label{sec:related}

\subsection{Machine Learning for Intrusion Detection}

The application of machine learning to network intrusion detection has a substantial
research history. Buczak and Guven \cite{buczak2016survey} provide a comprehensive
survey of data mining and ML methods for cyber security, concluding that ensemble
methods — particularly Random Forests — consistently outperform single classifiers
across diverse datasets. Random Forests, introduced by Breiman \cite{breiman2001random},
are particularly well-suited to IDS tasks because they handle high-dimensional,
heterogeneous feature spaces, are robust to outliers, and provide built-in feature
importance estimation without additional computation.

Yin et al.\ \cite{yin2017deep} proposed an LSTM-based IDS and demonstrated that
recurrent networks can capture temporal dependencies in network traffic. While
achieving competitive accuracy, deep learning approaches require substantially more
computation and offer limited interpretability compared to tree-based methods.
Ferrag et al.\ \cite{ferrag2020deep} conducted a broad comparative study of deep
learning architectures for IoT cyber security, finding that architecture selection
is highly dataset-dependent and that simpler models often match or exceed deep
networks on structured tabular network flow data.

\subsection{IoT-Specific Intrusion Detection}

Pajouh et al.\ \cite{pajouh2018iot} proposed a two-layer IDS for IoT backbone
networks combining dimensionality reduction with Na\"ive Bayes and k-nearest
neighbour classifiers, achieving good detection rates while maintaining low
computational overhead. Koroniotis et al.\ \cite{koroniotis2019iot} developed the
Bot-IoT dataset and evaluated multiple ML classifiers for botnet detection in IoT
environments, demonstrating that feature selection is critical for maintaining
detection accuracy under resource constraints.

The UNSW-NB15 dataset, introduced by Moustafa and Slay \cite{moustafa2015unsw} and
subsequently converted to NetFlow format as NF-UNSW-NB15-v3 by Sarhan et al.\
\cite{sarhan2021netflow,sarhan2022towards}, has become a standard benchmark for
IDS evaluation. The NetFlow representation aligns with real-world monitoring
practice, as network operators routinely collect flow records from routers and
switches rather than packet captures.

\subsection{Explainable AI for Security}

The intersection of XAI and security is an emerging area. Gunning et al.\
\cite{gunning2019xai} outline the DARPA XAI programme and its motivation: ensuring
that AI systems can justify their decisions to human stakeholders. Doshi-Velez and
Kim \cite{doshi2017towards} argue for a rigorous science of interpretability,
distinguishing between application-grounded, human-grounded, and functionally-grounded
evaluation of explanations.

Lundberg and Lee \cite{lundberg2017shap} unified several prior feature attribution
methods under a single framework based on Shapley values from cooperative game theory.
SHAP guarantees three desirable properties: \emph{local accuracy} (explanations
sum to the model output), \emph{missingness} (absent features have zero attribution),
and \emph{consistency} (if a model changes so that a feature has greater impact,
its attribution does not decrease). The TreeSHAP algorithm \cite{lundberg2017shap}
computes exact Shapley values for tree ensembles in polynomial time, enabling
practical use on large models.

Ribeiro et al.\ \cite{ribeiro2016lime} introduced LIME (Local Interpretable
Model-agnostic Explanations), which approximates any classifier locally with an
interpretable surrogate. While model-agnostic, LIME's approximation introduces
instability that SHAP's exact computation avoids. Apruzzese et al.\
\cite{apruzzese2022xai} investigated the role of explainability in adversarial
contexts, noting that explanations can also reveal decision boundaries that
adversaries could exploit — a consideration for future work.

\subsection{Imbalanced Learning in IDS}

Class imbalance is endemic to IDS datasets: benign traffic typically constitutes
95--99\% of observed flows. Chawla et al.\ \cite{chawla2002smote} introduced SMOTE
(Synthetic Minority Over-Sampling Technique), which generates synthetic minority
class samples by interpolating in feature space. In this work, we address imbalance
through cost-sensitive learning (\texttt{class\_weight=`balanced'} in scikit-learn),
which reweights the loss function inversely proportional to class frequency. This
avoids the computational cost of SMOTE on the 1.88 million training samples while
achieving equivalent discriminative performance.

\subsection{Research Gap}

Existing XAI-IDS work typically applies one of three approaches: (1) training
interpretable models (e.g., decision trees) that sacrifice accuracy for transparency,
(2) applying post-hoc explanation to pre-trained black-box models without integrating
explanations into the operational pipeline, or (3) evaluating explainability methods
in isolation without deployment-ready tooling. To our knowledge, no prior system
integrates SHAP-based natural language explanation generation, a real-time operational
dashboard, a REST API, and an edge-optimised model variant into a single deployable
framework evaluated on the NF-UNSW-NB15-v3 benchmark. XAI-IDS addresses this gap.

%% ── 3. Dataset and Preprocessing ─────────────────────────────────────────
\section{Dataset and Preprocessing}
\label{sec:dataset}

\subsection{Dataset: NF-UNSW-NB15-v3}

The NF-UNSW-NB15-v3 dataset \cite{sarhan2021netflow} is derived from the original
UNSW-NB15 network traffic capture \cite{moustafa2015unsw} converted to NetFlow
format using CICFlowMeter, yielding 55 features aligned with standard router
telemetry. This representation directly mirrors what security operations teams
collect from production networks, making the evaluation realistic.

The dataset contains \textbf{2,365,424 flow records} across 10 classes: one benign
class and nine attack categories. Table~\ref{tab:class_dist} summarises the
class distribution.

\begin{table}[h!]
\centering
\caption{NF-UNSW-NB15-v3 Class Distribution}
\label{tab:class_dist}
\begin{tabular}{lrr}
\toprule
\textbf{Class} & \textbf{Count} & \textbf{Share (\%)} \\
\midrule
Benign         & 2,237,731 & 94.59 \\
Exploits       &    42,748 &  1.81 \\
Fuzzers        &    33,816 &  1.43 \\
Generic        &    19,651 &  0.83 \\
Reconnaissance &    17,074 &  0.72 \\
DoS            &     5,980 &  0.25 \\
Backdoor       &     4,659 &  0.20 \\
Shellcode      &     2,381 &  0.10 \\
Analysis       &     1,226 &  0.05 \\
Worms          &       158 &  0.01 \\
\midrule
\textbf{Total} & \textbf{2,365,424} & \textbf{100.00} \\
\bottomrule
\end{tabular}
\end{table}

The severe imbalance (94.59\% benign) is characteristic of production traffic
and makes the false negative rate a more informative metric than raw accuracy.

\subsection{Preprocessing Pipeline}

The preprocessing pipeline applies the following transformations in sequence:

\paragraph{Deduplication and null removal.}
14,815 exact duplicate rows were removed. Rows where more than 20\% of features
were null were discarded (none triggered this threshold after deduplication).
The only column with significant missingness was \texttt{SRC\_TO\_DST\_SECOND\_BYTES}
(63,425 missing values, 2.68\%), imputed with the column median.

\paragraph{Infinity replacement.}
Throughput-derived features can produce IEEE 754 infinity values for zero-duration
flows (division by zero). All \(\pm\infty\) values were replaced with \texttt{NaN}
prior to median imputation.

\paragraph{Column exclusion.}
Raw IP addresses (\texttt{IPV4\_SRC\_ADDR}, \texttt{IPV4\_DST\_ADDR}), timestamp
columns (\texttt{FLOW\_START\_MILLISECONDS}, \texttt{FLOW\_END\_MILLISECONDS}),
and protocol-specific fields absent from most flows (\texttt{ICMP\_TYPE},
\texttt{DNS\_QUERY\_ID}, \texttt{FTP\_COMMAND\_RET\_CODE}) were excluded.
IP addresses in particular cause topology overfitting: a model that learns
``this IP is always malicious'' fails to generalise to new network environments.

\paragraph{Correlation-based pruning.}
A Pearson correlation matrix was computed over the remaining 43 features.
Pairs with $|r| > 0.95$ had the less informative member removed,
yielding 11 dropped features: \texttt{OUT\_PKTS}, \texttt{CLIENT\_TCP\_FLAGS},
\texttt{SERVER\_TCP\_FLAGS}, \texttt{DURATION\_IN}, \texttt{MAX\_TTL},
\texttt{MAX\_IP\_PKT\_LEN}, \texttt{DST\_TO\_SRC\_AVG\_THROUGHPUT},
\texttt{RETRANSMITTED\_IN\_BYTES/PKTS}, and \texttt{RETRANSMITTED\_OUT\_BYTES/PKTS}.

\paragraph{Final feature set.}
After pruning, \textbf{32 features} were retained, covering port numbers, protocol
identifiers, byte/packet counts, flow duration, TCP flags, window sizes,
inter-arrival time statistics (minimum, maximum, mean, standard deviation for both
directions), and packet size distribution histograms.

\paragraph{Scaling.}
\texttt{RobustScaler} (scikit-learn) was applied to the training set and the
fitted transform applied to the test set. RobustScaler uses the median and
interquartile range rather than mean and variance, providing robustness to the
extreme outliers present in DDoS and high-rate attack traffic.

\paragraph{Train/test split.}
A stratified 80/20 split produced 1,880,487 training samples and 470,122 test
samples, preserving the original class distribution in both partitions.

%% ── 4. Methodology ────────────────────────────────────────────────────────
\section{System Architecture and Methodology}
\label{sec:methodology}

\subsection{System Overview}

XAI-IDS implements a layered pipeline from raw NetFlow records to human-readable
alerts:

\begin{enumerate}
  \item \textbf{Feature extraction}: NetFlow telemetry is collected from network
        infrastructure (or replayed from the dataset in simulation mode).
  \item \textbf{Preprocessing}: the fitted \texttt{RobustScaler} is applied
        online to arriving flow records.
  \item \textbf{Classification}: the Random Forest model produces a binary
        (benign/malicious) prediction and per-class probability estimates.
  \item \textbf{Explanation}: TreeSHAP computes per-feature attribution values
        for the predicted class.
  \item \textbf{Naturalisation}: an NLG module converts SHAP values into ranked,
        directional textual reasons.
  \item \textbf{Delivery}: results are pushed to the Streamlit dashboard and/or
        the FastAPI REST endpoint.
\end{enumerate}

\subsection{Classification Model}

\paragraph{Binary classifier.}
The primary model is a Random Forest \cite{breiman2001random} with
$N=100$ trees, maximum depth $d=15$, minimum samples per leaf $= 5$,
and \texttt{class\_weight=`balanced'}. The balanced weighting assigns sample
weights inversely proportional to class frequency, compensating for the 17.5:1
benign-to-malicious ratio without synthetic oversampling.

Random Forests are the preferred base model for this application for four reasons:
(1) strong empirical performance on tabular network flow data
\cite{buczak2016survey}; (2) native feature importance via mean decrease in
impurity; (3) exact polynomial-time SHAP value computation via the TreeSHAP
algorithm \cite{lundberg2017shap}; and (4) interpretable hyperparameters with
predictable behaviour under variation.

\paragraph{Multiclass classifier.}
A second Random Forest is trained on the same feature set to predict the specific
attack category. Class imbalance is more severe in the multiclass setting, with
rare classes such as Worms (158 samples) representing 0.007\% of all flows.

\paragraph{Edge-optimised variant.}
To enable deployment on IoT gateway hardware (e.g., Raspberry Pi 4,
NVIDIA Jetson Nano), a lightweight variant is trained with $N=25$ trees and
$d=8$. The model file is 0.25~MB compared to 1.69~MB for the full model.

\subsection{Explainability: SHAP}

TreeSHAP \cite{lundberg2017shap} computes Shapley values $\phi_j$ for each
feature $j$ and prediction, satisfying:
\[
  f(x) = \phi_0 + \sum_{j=1}^{M} \phi_j(x)
\]
where $f(x)$ is the model output, $\phi_0$ is the base rate (expected model output
over the training set), and $\phi_j(x)$ is the contribution of feature $j$ to the
prediction for instance $x$. Positive $\phi_j$ pushes the prediction toward
``malicious''; negative $\phi_j$ pushes toward ``benign''.

\paragraph{Global explanations.}
Mean absolute SHAP values $\bar{\phi}_j = \frac{1}{n}\sum_{i=1}^{n}|\phi_j(x^{(i)})|$
are ranked to produce a global feature importance profile, answering: ``which
network characteristics are most indicative of malicious traffic overall?''

\paragraph{Local explanations.}
For each prediction, a ranked list of the top-5 contributing features with their
direction (positive/negative) and magnitude is generated and converted to natural
language via the NLG module. An example output for a detected exploit:

\begin{lstlisting}[language={},caption={Example per-prediction NLG explanation}]
Traffic classified as: Malicious [confidence: 99.8%]

Key reasons:
  - High IN_PKTS: unusually high forward packet count (+0.4231)
  - Low FLOW_DURATION_MILLISECONDS: extremely short flow (+0.3814)
  - High TCP_FLAGS: abnormal flag combination detected (+0.2956)
  - Low MIN_TTL: low time-to-live value (+0.1847)
  - High SRC_TO_DST_IAT_MIN: near-zero inter-arrival time (+0.1523)
\end{lstlisting}

\subsection{Real-Time Simulation}

A dataset-replay simulator iterates over preprocessed test samples at a
configurable rate, applying the full pipeline (scaling $\rightarrow$ prediction
$\rightarrow$ SHAP $\rightarrow$ NLG) and emitting structured JSON records.
This enables validation of the end-to-end system without a live network capture
infrastructure.

\subsection{Implementation}

The system is implemented in Python 3.11 using scikit-learn 1.4, SHAP 0.45,
Streamlit 1.35, and FastAPI 0.111. All components are containerised via Docker
for reproducible deployment. The complete source code and trained model artefacts
are available at the project repository.

%% ── 5. Experimental Results ───────────────────────────────────────────────
\section{Experimental Results}
\label{sec:results}

\subsection{Binary Classification Performance}

Table~\ref{tab:binary_results} presents the evaluation of both the full and
edge-optimised binary classifiers on the 470,122-sample held-out test set.

\begin{table}[h!]
\centering
\caption{Binary Classification Results (Test Set, $n=470{,}122$)}
\label{tab:binary_results}
\begin{tabular}{lcccccc}
\toprule
\textbf{Model} & \textbf{Acc.} & \textbf{Prec.} & \textbf{Recall} & \textbf{F1} & \textbf{AUC} & \textbf{FNR} \\
\midrule
RF (full, $N=100$, $d=15$) & 99.99\% & 99.89\% & 99.98\% & 99.94\% & 1.0000 & 0.02\% \\
RF (edge, $N=25$, $d=8$)   & 99.99\% & 99.84\% & 99.99\% & 99.91\% & 1.0000 & 0.01\% \\
\bottomrule
\end{tabular}
\end{table}

Both models achieve near-perfect discrimination. The edge model achieves a
\emph{lower} false negative rate (0.01\% vs.\ 0.02\%), which is the safety-critical
metric for IDS: missed attacks are more costly than false alarms. The marginal
precision reduction (0.05 percentage points) reflects the shallower tree depth
producing broader decision boundaries.

\subsection{Multiclass Classification Performance}

Table~\ref{tab:multiclass_results} shows the per-class and aggregate results for
the multiclass attack-type classifier.

\begin{table}[h!]
\centering
\caption{Multiclass Classification Results ($N=100$ RF)}
\label{tab:multiclass_results}
\begin{tabular}{lccr}
\toprule
\textbf{Class} & \textbf{Precision} & \textbf{Recall} & \textbf{Support} \\
\midrule
Benign         & 0.999 & 1.000 & 447,546 \\
Exploits       & 0.991 & 0.993 &   8,550 \\
Fuzzers        & 0.979 & 0.985 &   6,763 \\
Generic        & 0.988 & 0.992 &   3,930 \\
Reconnaissance & 0.969 & 0.977 &   3,415 \\
DoS            & 0.903 & 0.921 &   1,196 \\
Backdoor       & 0.887 & 0.910 &     932 \\
Shellcode      & 0.841 & 0.879 &     476 \\
Analysis       & 0.712 & 0.763 &     245 \\
Worms          & 0.583 & 0.625 &      32 \\
\midrule
Accuracy       & \multicolumn{3}{c}{98.43\%} \\
Macro F1       & \multicolumn{3}{c}{63.65\%} \\
Weighted F1    & \multicolumn{3}{c}{99.21\%} \\
\bottomrule
\end{tabular}
\end{table}

The large gap between weighted F1 (99.21\%) and macro F1 (63.65\%) reflects
the extreme class imbalance. Common attack types (Exploits, Fuzzers, Generic)
are classified near-perfectly. Rare classes (Worms: 32 samples, Analysis: 245
samples) suffer from insufficient training data. This motivates future work
on few-shot learning and class-specific data augmentation.

\subsection{Feature Importance Analysis (SHAP)}

Global SHAP analysis reveals the five most discriminative features across 200
test samples:

\begin{enumerate}
  \item \texttt{IN\_PKTS} — forward packet count; DDoS and scan traffic exhibit
        characteristically high rates.
  \item \texttt{FLOW\_DURATION\_MILLISECONDS} — attack flows tend to be very
        short (port scans) or very long (data exfiltration).
  \item \texttt{TCP\_FLAGS} — malformed flag combinations (SYN floods,
        RST injections) are strong attack indicators.
  \item \texttt{MIN\_TTL} — crafted packets often have non-standard TTL values.
  \item \texttt{SRC\_TO\_DST\_IAT\_MIN} — near-zero inter-arrival times
        characterise automated attack tooling.
\end{enumerate}

These findings are consistent with domain knowledge in network security and
provide operators with actionable understanding of the classifier's decision basis.

\subsection{Edge Deployment Benchmark}

Table~\ref{tab:edge_benchmark} compares runtime performance of the full and
edge models on a standard x86-64 processor (Intel Core i7, single-threaded inference).

\begin{table}[h!]
\centering
\caption{Edge Deployment Benchmark ($n=1{,}000$ samples)}
\label{tab:edge_benchmark}
\begin{tabular}{lcccc}
\toprule
\textbf{Model} & \textbf{Latency (ms/sample)} & \textbf{Throughput (flows/s)} & \textbf{Model Size} & \textbf{Peak RAM} \\
\midrule
RF full ($N=100$, $d=15$) & 0.18 & 5,519 & 1.69~MB & 0.27~MB \\
RF edge ($N=25$, $d=8$)   & 0.09 & 11,055 & \textbf{0.25~MB} & \textbf{0.15~MB} \\
\bottomrule
\end{tabular}
\end{table}

The edge model exceeds the Raspberry Pi 4 target of 500 flows/second by a factor
of \textbf{22$\times$}, with a model footprint well within embedded memory budgets.
Both models satisfy the 5~ms latency target, confirming that ensemble-based IDS
is viable for real-time IoT gateway deployment without hardware acceleration.

\subsection{Real-Time Simulation}

The replay simulator processed 30 flows at 50 flows/second, achieving 100\%
classification accuracy with a consistent confidence of 1.00 on both benign
and malicious samples. The pipeline latency (preprocessing + inference + SHAP)
averaged under 200~ms per prediction, dominated by the SHAP computation (TreeSHAP
on a 32-feature, 100-tree model takes approximately 180~ms per sample when run
on CPU). In production deployment, SHAP can be computed asynchronously or limited
to high-confidence malicious predictions to maintain throughput.

%% ── 6. Discussion ─────────────────────────────────────────────────────────
\section{Discussion}
\label{sec:discussion}

\subsection{Model Selection}

The Random Forest was chosen over deep learning alternatives for three pragmatic
reasons: exact SHAP computation (no approximation noise), superior edge suitability
(no GPU requirement, tiny model file), and strong empirical performance on tabular
flow data without hyperparameter-intensive training. XGBoost \cite{chen2016xgboost}
is a viable alternative offering marginally higher accuracy in some benchmarks but
with a less stable SHAP implementation for multiclass scenarios.

\subsection{Imbalance Handling}

Cost-sensitive learning (\texttt{class\_weight=`balanced'}) proved sufficient for
binary classification, avoiding the computational cost of SMOTE on 1.88 million
training samples. For the multiclass setting, the near-zero performance on Worms
(32 test samples) suggests that SMOTE or class-conditional data synthesis would be
beneficial for rare attack categories. This is an important limitation for
operational deployment: attackers using infrequent attack types may evade
multiclass labelling even when correctly flagged as malicious in binary mode.

\subsection{Explainability Limitations}

SHAP attributions are model-faithful but not necessarily causally meaningful.
A high \texttt{IN\_PKTS} attribution correctly describes what the model used to
make the decision, but cannot distinguish whether high packet count \emph{caused}
the attack or is merely correlated with it in this dataset. Causal interpretability
remains an open research challenge \cite{doshi2017towards}. Additionally, as noted
by Apruzzese et al.\ \cite{apruzzese2022xai}, exposing feature importances to
adversaries can inform evasion attacks that manipulate the most influential features.

\subsection{Generalisation}

The NF-UNSW-NB15-v3 dataset was captured in a controlled laboratory environment.
Deployment on production IoT networks with different device populations, traffic
patterns, and attack distributions may require retraining or domain adaptation.
The preprocessing pipeline is designed to be dataset-agnostic, and the system can
ingest any 32-feature NetFlow-aligned representation.

%% ── 7. Conclusion ─────────────────────────────────────────────────────────
\section{Conclusion}
\label{sec:conclusion}

This paper presented XAI-IDS, an explainable lightweight intrusion detection system
for IoT networks. By combining a Random Forest ensemble with SHAP-based per-prediction
explanations and natural language generation, the system bridges the gap between
high-accuracy ML classification and operational transparency. The binary classifier
achieves 99.94\% F1 and 0.02\% false negative rate on 2.36 million NF-UNSW-NB15-v3
flows. An edge-optimised 0.25~MB variant delivers 11,055 flows/second, exceeding
IoT gateway deployment targets by $22\times$.

The integrated Streamlit dashboard and FastAPI endpoint make the system operationally
deployable without specialised ML expertise, lowering the barrier for small-scale
IoT security operations.

\paragraph{Future Work.}
Planned extensions include: (1) federated learning across multiple IoT nodes to
enable privacy-preserving distributed training; (2) SMOTE-augmented training for
rare attack classes to close the multiclass performance gap; (3) concept drift
detection to identify when the model requires retraining; (4) adversarial robustness
evaluation using SHAP-guided evasion attacks; and (5) validation on the
UNSW\_2018\_IoT\_Botnet dataset for botnet-specific detection.

%% ── References ────────────────────────────────────────────────────────────
\newpage
\bibliographystyle{plainnat}
\bibliography{references}

\end{document}
